\documentclass{article}
\usepackage[utf8]{inputenc}
\usepackage[T1]{fontenc}
\usepackage[margin=2cm]{geometry}
\usepackage{hyperref}
\sloppy
\begin{document}
\section*{Historical source}
Encyclopaedia Britannica scan, source row 2001. Parse 3.5 extracted the text. Super 3.5 VL corrected it from an enlarged crop of each block. It can still contain recognition errors.
Page headers, footers, and page numbers are removed.

\subsection*{Untitled continuation (E0172)}
\textit{Continues an entry from an earlier page.}

bers of black cattle and horses are brought hither from Poland and Hungary for sale, those bred in the country not being sufficient; but of sheep, goats, game, and venison, they have great plenty. As for wild beasts, here are lynxes, foxes, weasels, otters, and beavers. The rivers, lakes, and ponds, yield fish of several sorts, particularly sturgeons several ells in length, and salmon. Besides a number of smaller streams to water this country, there is the Oder, which traverses it almost from one end to the other; and the Vistula, which after a pretty long course through it enters Poland. The number of the cities and market-towns is said to be about 200, the county of Glatz included, and that of the villages 5000. The inhabitants who are computed to be about 1,821,065 are a mixture of Germans, Poles, and Moravians. The language generally spoken is German; but in some places the vulgar tongue is a dialect of the Slavonic. The states consist of the princes and dukes, and those called plate-lords, with the nobility, who are immediately subject to the sovereign, and the representatives of the chief cities; but since the country fell under the dominion of the king of Prussia, no diets have been held. The king, however, when he took possession of the country, confirmed all the other privileges of the inhabitants. With respect to religion, not only Protestants, but Papists, Jews, and Greeks, enjoy full liberty of conscience. The greatest part of Silesia lies in the diocese of Breslaw, but some part of it in the Polish dioceses of Posen and Cracow. The bishop of Breslaw stands immediately under the pope with regard to spirituals; but all ecclesiastical benefices, not excepting the fee of Breslaw, is in the king's gift. Besides Latin schools, colleges, and seminaries, at Breslaw is an university, and at Lignitz an academy for martial exercises. The principal manufactures here are woollens, linens, and cottons of several sorts, with hats, glass-ware, gunpowder, and iron manufactures. Of these there is a considerable exportation. Accounts are generally kept in rix-dollars, silver groschens, and ducats. With respect to its revolutions and present government, it was long a part of the kingdom of Poland; afterwards it had several dukes and petty princes for its sovereigns, who by degrees became subject to the kings of Bohemia, until at last King Charles IV. incorporated the whole duchy with Bohemia; and thus it continued in the possession of the house of Austria, until the king of Prussia in 1742, taking advantage of the troubles that ensued upon the death of the emperor Charles VI. and pretending a kind of claim, wrested a great part of it, together with the county of Glatz, from his daughter and heirs Maria Theresa, the late empress dowager; so that now only a small part of it is possessed by the house of Austria, and connected with the empire, the rest being governed by the king of Prussia, without acknowledging any sort of dependence on the crown of Bohemia or the empire. For the administration of justice in all civil, criminal, and feudal cases, and such as relate to the revenue, the king of Prussia has established three supreme judicatories, to which an appeal lies from all the inferior ones, and from which, when the sum exceeds 500 rix-dollars, causes may be moved to Berlin. The Lutheran churches and schools are under the inspection of the upper consistories, and those of the Papists under that of the bishop's court at Breslaw; but from both an appeal lies to the tribunal at Berlin.

As to the revenue, the excise here is levied only in the walled towns, being on the same footing as in the marquisate of Brandenburg; but in the rest of the country the contributions are fixed, and the same both in peace and war. The several branches of the revenue are under the management of the war and domain offices of Breilaw and Glogau. The whole revenue arising to the king of Prussia from Silesia and the county of Glatz amounts to about four millions of rix-dollars per annum.

Silesia is divided into Upper and Lower, and each of these again into principalities and lordships; of some of which both the property and jurisdiction belong immediately to the foreigner, but of others to his subjects and vassals. In regard to the character of the people, the boors are accounted very dull and stupid; but of those of a higher rank, many have distinguished themselves by their wit and learning, as well as by their military and political talents.

\subsection*{SILESIAN EARTH (E0173)}
SILESIAN EARTH, in the Materia Medica, a fine astringent bole. It is very heavy, of a firm compact texture, and in colour of a brownish yellow. It breaks easily between the fingers, and does not stain the hands; is naturally of a smooth surface, is readily diffusible in water, and melts freely into a butter-like fulgence in the mouth. It leaves no grittiness between the teeth, and does not ferment with acids. It is found in the perpendicular fissures of rocks near the gold mines in Hungary.

\subsection*{SILICERNIUM (E0174)}
SILICERNIUM, among the Romans, was a feast of a private nature, provided for the dead some time after the funeral. It consisted of beans, lettuces, bread, eggs, \&c. These were laid upon the tomb, and they foolishly believed that the dead would come out for the repast. What was left was generally burnt on the stone. The word silicernium is derived from silex and cera, i. e. "a supper upon a stone." Eating what had thus been provided for the dead, was esteemed a mark of the most miserable poverty. A similar entertainment was made by the Greeks at the tombs of the deceased; but it was usual among them to treat the ghouls with the fragments from the feast of the living. See FUNERAL AND INFERIUS.

\subsection*{SILEX (E0175)}
SILEX. See FLINT.

\subsection*{SILICEOUS EARTHS (E0176)}
SILICEOUS EARTHS. See SILICA, CHEMISTRY Index.

\subsection*{SILIUS ITALICUS, CAIUS (E0177)}
SILIUS ITALICUS, CAIUS, an ancient Roman poet, and author of an epic poem in 17 books, which contains an history of the second Punic war, so famous for having decided the empire of the world in savour of the Romans. He was born in the reign of Tiberius, and is supposed to have derived the name of Italicus from the place of his birth; but whether he was born at Italica in Spain, or at Corfinium in Italy, which, according to Strabo, had the name of Italica given it during the Social war, is a point which cannot be known: though, if his birth had happened at either of those places, the grammarians would tell us, that he should have been called Italicensis, and not Italicus. When he came to Rome, he applied himself to the bar; and, by a close imitation of Cicero, succeeded so well, that he became a celebrated advocate and most accomplished orator. His merit and character recommended him to the highest offices in the republic, even to the consulship, of which he was possessed when Nero died. He is said to


\section*{Source record}
\begin{description}
\item[Scan ID] \texttt{britannica:309f765b1ecc83e0a7f7762e44b4a927a7555443:train:2001}
\item[Image file] \texttt{page-2001-93040a5ab362.jpg}
\item[Image SHA-256] \texttt{93040a5ab3625efc0259c9878caea3e2285bc73aae0c1fda452f9f5198dedd67}
\item[Dataset] NationalLibraryOfScotland/encyclopaedia\_britannica\_illustrated at revision \texttt{309f765b1ecc83e0a7f7762e44b4a927a7555443} (CC0-1.0)
\item[Extraction] model \texttt{nvidia/nemotron-3.5-parse}, raw reply \texttt{extract/raw/page-2001-93040a5ab362.json}
\item[Raw reply SHA-256] \texttt{4077eea4176c84367ec7d0e6a46bcc0112e023636a15ac13b8504e03cc96af91}
\item[Correction] model \texttt{nvidia/Nemotron-3.5-Super-VL}, block crops, raw replies \texttt{correct/raw/crop\_nothink/page-2001-93040a5ab362-b*.json}; parse: 1, super\_vl: 8
\item[Running header] S I L [ 355 ]
\end{description}
\end{document}
